\documentclass[11pt]{article}
\usepackage[T1]{fontenc}
\usepackage[utf8]{inputenc}
\usepackage[english]{babel}
\usepackage{amsmath,amssymb,mathtools,array,booktabs}
\usepackage{geometry}
\usepackage{microtype}
\geometry{a4paper,margin=1.45cm}
\setlength{\parindent}{1.25em}
\setlength{\parskip}{0pt}
\makeatletter
\renewcommand\footnotesize{\@setfontsize\footnotesize{7.5}{8.5}\selectfont}
\makeatother
\allowdisplaybreaks
\setlength{\emergencystretch}{12em}
\sloppy
\hbadness=10000
\vbadness=10000
\hfuzz=2pt
\vfuzz=10pt
\raggedbottom
\begin{document}

\begin{center}
{\Large\bfseries 9. The Most General Domains from Integral Transcendental Numbers}\par
\vspace{1.2em}
Math. Ann. 77 (1916), pp. 103--128
\end{center}

By means of the well-ordering theorem, E. Zermelo constructed a domain of
``integral transcendental numbers,'' that is, a numerical domain $\mathfrak G$
having the following abstract properties:
\begin{enumerate}
\item[I.] The sum, difference, and product of two numbers from $\mathfrak G$ is
again a number from $\mathfrak G$.
\item[II.] Every real or complex number is the quotient of two numbers from
$\mathfrak G$.
\item[III.] Every integral rational (respectively algebraic) number is an element
of $\mathfrak G$.
\item[IV.] No non-integral rational (respectively algebraic) number belongs to
$\mathfrak G$.\footnote{E. Zermelo, Über ganze transzendente Zahlen,
Math. Ann. 75, p. 434 (1914).}
\end{enumerate}

The construction rests on the fact that the well-ordering theorem implies the
existence of an \emph{algebraic basis} of all numbers: namely, a \emph{system}
$H$ \emph{of numbers} $\eta$ among which no algebraic relations exist, but by
means of which all remaining numbers can be expressed algebraically. Because of
their algebraic independence, these basis numbers $\eta$ can be regarded as
indeterminates, and the desired domain therefore becomes an integral domain of
rational and algebraic functions of indeterminates $\eta$,\footnote{For this
reason the considerations of the present paper partly run parallel to those of the
author's earlier paper: Körper und Systeme rationaler Funktionen, Math. Ann. 76,
p. 161 (1915).} whose coefficients belong to the field $K$ of all algebraic
numbers. The special domain $\mathfrak G_\eta$ constructed by Zermelo is then
characterized by the following basis properties:

$\mathfrak G_\eta$ contains:
\begin{enumerate}
\item[1)] all basis numbers $\eta$,
\item[2)] all polynomials in the $\eta$ with integral\footnote{Throughout,
``integral'' is to mean that the coefficients belong to the integral domain $[K]$
of all algebraic integers. For the definition of $\mathfrak G_\eta$, it would
also suffice in 1)--5) to consider only rational integer coefficients; for more
general domains, however, the basis properties would then become less simple.}
coefficients,
\item[3)] all integral algebraic functions of the $\eta$ with integral
coefficients.
\end{enumerate}
$\mathfrak G_\eta$ excludes:
\begin{enumerate}
\item[4)] all rationally fractional functions of the $\eta$ with integral
coefficients,
\item[5)] all algebraically fractional functions of the $\eta$ with integral
coefficients.\footnote{Zermelo, loc. cit., § 3.}
\end{enumerate}

It is now easy to see (compare the examples in §§ 2 and 3) that there are
domains $\mathfrak G$ essentially different from the Zermelo domain
$\mathfrak G_\eta$: that is, domains $\mathfrak G$ for which, under no
well-ordering, do all basis properties 1)--5) hold, and which consequently cannot
be generated by Zermelo's construction under any well-ordering. In other words,
there are domains $\mathfrak G$ that cannot be mapped one-to-one and
isomorphically onto $\mathfrak G_\eta$. Thus the question arises of the most
general domains from integral transcendental numbers; it will be answered
completely in what follows.

For this purpose one must first determine which of the basis properties are
consequences of the abstract defining properties and which are caused by the
special construction. It is shown (§ 1) that for every prescribed domain
$\mathfrak G$ there is a well-ordering such that basis properties 1) and 2) are
satisfied; hence every domain $\mathfrak G$ can be constructed by means of an
algebraic basis of all numbers. None of the further basis properties need be
satisfied (§§ 2 and 3). In particular it follows that a further abstract property
of the Zermelo domain does not belong to all domains $\mathfrak G$: algebraic
integral closedness. It can be formulated as follows:
\begin{enumerate}
\item[V.] Every number that depends algebraically-integrally and integrally on
finitely many numbers from $\mathfrak G$ belongs to $\mathfrak G$.
\end{enumerate}

The special domains for which V holds in addition to I--IV will be called domains
$\mathfrak H$ from algebraically-integral transcendental numbers. For the most
general domains $\mathfrak H$ one obtains the simple results (§ 4):
\begin{enumerate}
\item[a)] The intersection of all domains $\mathfrak H$ constructible by means of
the same basis $H$ is given by all integral algebraic functions of the basis, that
is, by the Zermelo domain $\mathfrak G_\eta$, which is itself a domain
$\mathfrak H$. (This also gives an abstract characterization of the Zermelo
domain.)
\item[b)] Every domain $\mathfrak H$ arises by algebraically-integral
adjunction\footnote{Explanation of the expression in § 4.} of an arbitrary system
$\mathfrak S(\eta)$ to $\mathfrak G_\eta$, where $\mathfrak S(\eta)$ has to
satisfy only the single condition that no algebraically fractional number enters
the domain through the adjunction.
\end{enumerate}

The results for the most general domains $\mathfrak G$, when an algebraic basis
is taken as foundation, are less simple. Here the intersection of all domains
$\mathfrak G$ is itself not a domain $\mathfrak G$, and the construction of the
most general domains is consequently harder to survey (§ 5).

In order to obtain results analogous to those for the domains $\mathfrak H$, the
algebraic basis is replaced by a \emph{rational basis} $\Theta$, that is, by a
\emph{system} $\Theta$ \emph{of numbers} $\vartheta$ which permits all numbers
to be expressed rationally, while under at least one well-ordering no basis number
can be expressed rationally by finitely many preceding ones. This rational basis
$\Theta$ must still be subject to the side condition that the integral domain
$[\Theta]$ consisting of all integral polynomials in the $\vartheta$ contains no
algebraically fractional number. (The construction of such a basis with the side
condition is shown in § 7.) Then for the most general domains $\mathfrak G$ ---
which one could also call domains from rationally integral transcendental numbers
--- one again obtains the simple results (§ 6):
\begin{enumerate}
\item[a)] The intersection of all domains $\mathfrak G$ constructible by means of
the same rational basis $\Theta$ is given by all integral rational functions of the
basis, that is, by the domain $[\Theta]$, which is itself a domain $\mathfrak G$.
\item[b)] Every domain $\mathfrak G$ arises by rationally integral adjunction of
an arbitrary system $\mathfrak S(\vartheta)$ to $[\Theta]$, where
$\mathfrak S(\vartheta)$ has to satisfy only the single condition that no
algebraically fractional number enters the domain through the adjunction.
\end{enumerate}

To have complete analogy with the domains $\mathfrak H$, one would have to omit
the algebraic numbers from the defining properties III and IV for the domains
$\mathfrak G$. With this modification of III and IV, the construction of the
integral elements by means of a rational basis can be transferred to an arbitrary
field (§ 10), whereas Zermelo's construction presupposes the algebraic closedness
of the field.

As $H$ runs through every algebraic basis, respectively as $\Theta$ runs through
every rational basis satisfying the side condition, the results a) and b) give the
totality of all domains $\mathfrak H$ and $\mathfrak G$. If all isomorphic
domains are collected into a class, then from this totality one can select a subset
which represents the totality of all classes --- at least countably infinitely many
by § 9 (§ 8).

\subsection*{§ 1. Proof of basis properties 1) and 2) for all domains
$\mathfrak G$.}

Let $\mathfrak G$ be any prescribed domain from integral transcendental numbers,
so that it has the abstract properties I--IV. Following the procedure applied by
E. Zermelo to the field of all complex numbers, we select from $\mathfrak G$ an
\emph{algebraic basis} $H$ \emph{of all numbers from $\mathfrak G$}. Since by
II every real or complex number can be represented as the quotient of two numbers
from $\mathfrak G$, $H$ is at the same time an algebraic basis of \emph{all}
numbers. Thus, for every prescribed domain $\mathfrak G$, the algebraic basis of
all numbers can be chosen so that it belongs to $\mathfrak G$; hence basis
property 1) is satisfied. Moreover, by III, $\mathfrak G$ contains the integral
domain $[K]$ of all algebraic integers, and therefore, by I, also all polynomials
in the basis elements $\eta$ with coefficients from $[K]$, that is, all integral
polynomials in the $\eta$, which form the integral domain $[H]$. Hence basis
property 2) is also always satisfied; in other words, every domain $\mathfrak G$
can be constructed by means of an algebraic basis $H$ of all numbers in such a way
that $\mathfrak G$ contains the integral domain $[H]$ as a divisor.

\textbf{Remark.} Nevertheless, starting from a basis $H$, one can of course
construct a domain $\mathfrak G$ for which basis properties 1) and 2) do not hold
with respect to $H$ itself, but do hold with respect to some other basis $H'$.
For example, let $\mathfrak G'$ arise from the Zermelo domain
$\mathfrak G_\eta$ by replacing, throughout the whole construction, some basis
element $\eta$ by a polynomial $f(\eta)$. Then $\mathfrak G'$ contains neither
$\eta$ nor $[H]$; but if in the basis $H$ the basis element $\eta$ is replaced by
$\xi=f(\eta)$ --- whereby $H$ again goes over into an algebraic basis $H'$ ---,
then basis properties 1) and 2) are satisfied for $\mathfrak G'$ with respect to
$H'$.

\subsection*{§ 2. Exclusion of basis properties 4) and 5).}

We now show that the domains $\mathfrak G$ need not satisfy any of the further
basis properties. First 4) and 5) are excluded by replacing, in Zermelo's
construction, the integral domain $[H]$ by a domain of integral rational
``functionals.''

According to Weber,\footnote{H. Weber, Lehrbuch der Algebra. Kleine Ausgabe
§§ 96 and 97.} a \emph{rational functional} means every rational function with
rational number coefficients in the following uniquely determined representation:
\[
 A(\eta_1\cdots \eta_t)
   = a\cdot\frac{E_1(\eta_1\cdots \eta_t)}{E_2(\eta_1\cdots \eta_t)},
\]
where $E_1,E_2$ are relatively prime primitive polynomials with rational integer
coefficients, and $a$ is a positive rational number (the absolute value of $A$).
If, in particular, $a$ is a \emph{rational integer}, then $A$ is called an
\emph{integral rational functional}. As special cases, the integral rational
functionals include the rational integers and the polynomials with rational
integer coefficients, while rationally fractional numbers and polynomials with
rationally fractional coefficients are to be counted among the fractional rational
functionals.

Let the domain $\mathfrak G$ now consist of the totality $\mathfrak F$ of all
integral rational functionals of finitely many basis elements $\eta$ at a time,
and of all functions which depend algebraically-integrally on $\mathfrak F$ ---
that is, which satisfy some equation whose leading coefficient is unity and whose
remaining coefficients are integral rational functionals.

The proof of properties I--IV for $\mathfrak G$ runs exactly parallel to the
corresponding proof for Zermelo and will only be indicated briefly. First, II and
III are clear, since $\mathfrak G$ contains the Zermelo domain
$\mathfrak G_\eta$ as a divisor. By Gauss's theorem on polynomials with rational
integer coefficients, sums, differences, and products of integral rational
functionals are again integral rational functionals; hence $\mathfrak F$ forms an
integral domain, and by familiar arguments\footnote{Compare, for example, Weber,
§ 97, 8.} $\mathfrak G$ also becomes an integral domain; thus I is fulfilled.
Further, from the extended Gauss theorem\footnote{Weber, § 20, 5.} follow the two
auxiliary propositions: ``If $z$ depends algebraically-integrally on
$\mathfrak F$ by means of \emph{some} equation, then it also does so by means of
the \emph{irreducible} equation which $z$ satisfies with respect to the field of
all rational functionals,''\footnote{Compare Weber, § 97, 4.} and ``the equation
irreducible over the field of all rational numbers and satisfied by an algebraic
number is also irreducible over the field of all rational functionals.'' Thus
$\mathfrak G$ can contain no fractional algebraic number; hence IV, and therefore
all abstract properties I--IV, have been verified.

On the other hand, under no well-ordering can a basis be selected from
$\mathfrak G$ in such a way that $\mathfrak G$ excludes all rationally and
algebraically fractional functions of the basis elements. For let $z(\eta)$ be any
function of the domain to be chosen as a basis element, hence defined by an
equation
\[
 z^k+A_1(\eta)z^{k-1}+\cdots+A_k(\eta)=0,
\]
where
\[
 A_i(\eta)=a_i\cdot\frac{E_1(\eta)}{E_2(\eta)}
\]
are integral rational functionals. Then the function
$\frac{a_k}{z}$ belongs to $\mathfrak G$, and likewise
$\frac{a_k}{\sqrt z}$, $\frac{a_k}{\sqrt[3]z}$, etc.\footnote{Quite analogously,
every algebraic function of the $\eta$ can be transformed into an element of the
functional domain $\mathfrak G$ by multiplication by a rational integer. Thus
this domain has the property that every complex number can be represented as the
quotient of an element from $\mathfrak G$ and a rational integer, in analogy with
the algebraic numbers.} Hence basis properties 4) and 5) are excluded for the
totality of the domains $\mathfrak G$.

\textbf{Remark.} An example in § 3 will show that $\mathfrak G$ may, under every
well-ordering, contain rationally fractional functionals without containing a
rationally or algebraically fractional number.

\subsection*{§ 3. Exclusion of basis property 3).}

For excluding basis property 3), we use a generalized \emph{module domain} in
which the arguments of the polynomials of the domain are no longer the
indeterminates, but all integral algebraic functions of the indeterminates,
\footnote{By ``integral algebraic functions'' are always meant functions with
\emph{integral} coefficients; that is, functions satisfying some equation whose
leading coefficient is unity and whose remaining coefficients are polynomials in
the $\eta$ with algebraic-integer coefficients --- polynomials from $[H]$. In
particular, all polynomials from $[H]$ belong to the integral algebraic
functions.} while the indeterminates themselves take over the role of a module
basis.

Let $\mathfrak G$ consist of \emph{all elements}
\begin{equation}
 z=\alpha+g_1(\eta)\eta_1+g_2(\eta)\eta_2+\cdots+g_t(\eta)\eta_t,
\tag{1}
\end{equation}
\emph{where $\alpha$ runs through all algebraic integers,
$\eta_1\cdots\eta_t$ through all basis elements, and for each fixed combination
$\eta_1\cdots\eta_t$, the functions
$g_1(\eta)\cdots g_t(\eta)$ individually run through all integral algebraic
functions of the $\eta$.}\footnote{The module property belongs to the elements
$z-\alpha$.}

It is easy to verify that properties I--IV hold for this module domain. First I
holds, since the sum, difference, and product of two elements $z$ again takes the
form (1). The domain $\mathfrak G$ also contains all basis elements $\eta$, and
besides all elements $\eta\cdot g(\eta)$, where $g(\eta)$ runs through all
integral algebraic functions of the $\eta$. Hence all integral algebraic functions
of the $\eta$, and therefore all algebraic functions of the $\eta$ altogether,
can be represented as quotients of two elements of $\mathfrak G$, proving II.
In the representation (1), all algebraic integers are also included in particular
--- for $g_1=0\cdots g_t=0$ ---; but $z$ cannot be equal to a fractional
algebraic number. For if $z$ represents an algebraic number $\beta$, then from
\[
 \beta=\alpha+g_1(\eta)\eta_1+\cdots+g_t(\eta)\eta_t
\]
identically in the $\eta$, necessarily $\beta=\alpha$. This is shown by the
specialization
\[
 \eta_1=0\cdots \eta_t=0,
\]
under which the multipliers $g_1(\eta)\cdots g_t(\eta)$ remain finite as integral
algebraic functions, since they go over into integral algebraic functions or
numbers.\footnote{This more detailed proof of IV is given because of the expanded
example at the end of the paragraph. Here the remark would suffice that, by
construction, $z$ is an integral algebraic function.} Thus conditions I--IV are
satisfied for $\mathfrak G$.

On the other hand, under no well-ordering can a basis be selected from
$\mathfrak G$ in such a way that all integral algebraic functions of the basis
elements belong to $\mathfrak G$. For let $z(\eta)$ be any function of the domain
to be chosen as a basis element --- which therefore must actually contain the
indeterminates $\eta$ ---. We shall show that, from a certain finite value of
$\nu$ onward, depending on $z$, the functions
\[
 \xi=\sqrt[\nu]{z-\alpha}
\]
can no longer belong to $\mathfrak G$.

Indeed, suppose that $\xi$ belongs to $\mathfrak G$. Then, by (1), it is of the
form
\[
 \xi=\gamma+h_1(\eta)\eta_{i_1}+h_2(\eta)\eta_{i_2}
      +\cdots+h_\tau(\eta)\eta_{i_\tau},
\]
where again $h_1(\eta)\cdots h_\tau(\eta)$ are integral algebraic functions of
the $\eta$, and $\eta_{i_1}\cdots\eta_{i_\tau}$ are arbitrary basis elements,
which in particular may also coincide with $\eta_1\cdots\eta_t$. First it must be
shown that the algebraic integer $\gamma$ is zero. Since
$h_1(\eta)\cdots h_\tau(\eta)$ are integral algebraic functions, $\xi$ becomes
equal to $\gamma$ for
$\eta_{i_1}=0\cdots\eta_{i_\tau}=0$ with arbitrary values of the remaining
$\eta$; in particular, then, $\xi$ is equal to $\gamma$ for
\[
 \eta_{i_1}=0\cdots \eta_{i_\tau}=0; \qquad
 \eta_1=0\cdots \eta_t=0.
\]
But under this specialization the function $(z-\alpha)$ vanishes by (1), and
therefore so does $\xi$; hence $\gamma=0$, and one has
\[
 \xi=h_1(\eta)\eta_{i_1}+h_2(\eta)\eta_{i_2}
      +\cdots+h_\tau(\eta)\eta_{i_\tau}.
\]
Raising this to the $\nu$-th power gives
\[
 \xi^\nu=z-\alpha=F_\nu(\eta),
\]
where $F_\nu(\eta)$ denotes a homogeneous, not identically vanishing, form of
$\nu$-th dimension in $\eta_{i_1}\cdots\eta_{i_\tau}$, whose coefficients are
integral algebraic functions of the $\eta$. Now $z-\alpha$, which by (1) is an
integral algebraic function of the $\eta$, satisfies an equation irreducible with
respect to $[H]$:
\[
 (z-\alpha)^\chi+B(\eta)(z-\alpha)^{\chi-1}+\cdots+C(\eta)=0,
\]
where the last coefficient $C(\eta)$ --- the product of $(z-\alpha)$ with its
conjugates --- is a polynomial; let its degree be $\lambda$. On the other hand,
from $z-\alpha=F_\nu(\eta)$, by multiplying $F_\nu(\eta)$ by the corresponding
conjugates, one obtains
\[
 C(\eta)=G_{\nu\chi}(\eta),
\]
where $G_{\nu\chi}(\eta)$ denotes a homogeneous form of $\nu\chi$-th dimension
in $\eta_{i_1}\cdots\eta_{i_\tau}$, whose coefficients are polynomials in the
$\eta$. Thus $G_{\nu\chi}$ is at least of degree $\nu\chi$ in $\eta$, or
$\lambda\geq\nu\chi$; that is, the functions $\sqrt[\nu]{z-\alpha}$ for which
$\nu>\lambda/\chi$ do not belong to $\mathfrak G$. Basis property 3) is therefore
excluded for the totality of the domains $\mathfrak G$.

\textbf{Remark.} A slight generalization of the domain just constructed leads to
a domain $\mathfrak G$ which, under every well-ordering, also contains fractional
functionals. Let $\mathfrak G$ consist of all elements
\[
 z=\alpha+\gamma_1(\eta)\eta_1+\gamma_2(\eta)\eta_2
      +\cdots+\gamma_t(\eta)\eta_t,
\]
where now $\gamma(\eta)$ runs through all algebraically integral, but not
necessarily integral, functions of the $\eta$. The proof of properties I--IV
remains the same as above. But since $\mathfrak G$, besides every function $z$,
also contains $a\cdot(z-\alpha)$, where $a$ denotes any rational or
algebraically fractional number, fractional functionals also occur under every
well-ordering.

\subsection*{§ 4. The most general domains from algebraically-integral
transcendental numbers.}

For a more convenient formulation of what follows, introduce the expression:
``an element $z$ depends algebraically-integrally on $\mathfrak T$,'' meaning
that $z$ satisfies some equation whose leading coefficient is unity and whose
remaining coefficients are integral polynomials in the elements of $\mathfrak T$,
where $\mathfrak T$ denotes any prescribed domain.\footnote{For an integral
domain $\mathfrak T$, the leading coefficient is unity and the remaining
coefficients are elements of $\mathfrak T$.}

A domain $\mathfrak T$ is called \emph{algebraically-integrally closed} if it
satisfies the condition (compare the introduction):
\begin{enumerate}
\item[V.] Every element that depends algebraically-integrally on $\mathfrak T$
belongs to $\mathfrak T$.
\end{enumerate}

We first show that the abstract property V belongs to the Zermelo domain
$\mathfrak G_\eta$. Let an element $z$ depend algebraically-integrally on
$\mathfrak G_\eta$ by means of an equation $f(z)=0$. Then multiplication of
$f(z)$ by its conjugates with respect to $[H]$ shows that $z$ also depends
algebraically-integrally on $[H]$, and hence belongs to $\mathfrak G_\eta$.
The algebraically-integral closedness of the functional domain considered in § 2
is shown in the same way.

That property V does not belong to every domain $\mathfrak G$ is shown by the
module domain considered in § 3, which indeed does not possess basis property 3).
More precisely, § 3 showed that the module domain is algebraically-integrally
closed with respect to no transcendental element of the domain, whereas, by III
and IV, this is of course the case with respect to every algebraic element of the
domain.

This leads us first to select, from the totality of domains $\mathfrak G$, those
which also possess the abstract property V; these are to be called ``domains
$\mathfrak H$ from algebraically-integral transcendental numbers.''

Let $\mathfrak H$ be such a domain, and let $H$ be an algebraic basis of all
numbers from $\mathfrak H$; by II, $H$ is at the same time an algebraic basis of
all numbers. By III, I, and V, $\mathfrak H$ contains, besides $H$, all integral
algebraic functions of the $\eta$, that is, the Zermelo domain
$\mathfrak G_\eta$. Thus, in analogy with § 1: every domain $\mathfrak H$ from
algebraically-integral transcendental numbers can be constructed by means of an
algebraic basis of all numbers in such a way that $\mathfrak H$ contains the
Zermelo domain $\mathfrak G_\eta$ as a divisor.

Now let $H$ be any algebraic basis of all numbers, and let $\mathfrak M(H)$
denote the totality of all domains $\mathfrak H$ containing $H$. Every domain
$\mathfrak H$ in $\mathfrak M(H)$ contains, as just proved, the Zermelo domain
$\mathfrak G_\eta$ as a divisor; and since $\mathfrak G_\eta$ itself is a domain
$\mathfrak H$, $\mathfrak G_\eta$ is the greatest common divisor, or
intersection, of all domains $\mathfrak H$ from $\mathfrak M(H)$. Thus
$\mathfrak G_\eta$ has been defined abstractly. It also follows directly that
$\mathfrak M(H)$ is closed under intersection: that is, the intersection of all
domains $\mathfrak H$ of any subsystem of $\mathfrak M(H)$ again belongs to
$\mathfrak M(H)$. For, as an intersection, I and V are satisfied; II and III
follow because $\mathfrak G_\eta$ is a divisor; IV follows because the
intersection is a divisor of $\mathfrak H$.

As the example of the functional domain in § 2 shows, the domains $\mathfrak H$
in general contain further functions besides $\mathfrak G_\eta$. Let
$\psi(\eta)$ be such a rational or algebraic function of the $\eta$. Then, by I,
$\mathfrak H$ also contains all integral rational combinations of $\psi$ and
finitely many functions from $\mathfrak G_\eta$ at a time. The integral domain so
arising --- consisting of all polynomials in $\psi$ with coefficients from
$\mathfrak G_\eta$ --- will be denoted by $[\mathfrak G_\eta,\psi]$, and the
process leading to it will be called ``rationally integral adjunction of $\psi$
to $\mathfrak G_\eta$.'' By V, $\mathfrak H$ further contains all functions which
depend algebraically-integrally on $[\mathfrak G_\eta,\psi]$; the domain so
arising will be denoted by $\{\mathfrak G_\eta,\psi\}$, and the process leading
to it by ``algebraically integral adjunction of $\psi$ to $\mathfrak G_\eta$.''
The domain $\{\mathfrak G_\eta,\psi\}$ consists of all integral algebraic
functions of $\psi$ with coefficients from $\mathfrak G_\eta$, and is therefore
an integral domain by familiar arguments.\footnote{Compare, for example, Weber,
§ 97, 6 and 7.}

We show that $\{\mathfrak G_\eta,\psi\}$ is also algebraically-integrally closed,
that is, satisfies condition V. Indeed, let $z$ depend algebraically-integrally
on $\{\mathfrak G_\eta,\psi\}$ by means of an equation
\[
 f(z;a\cdots c)=z^\nu+a z^{\nu-1}+\cdots+c=0,
\]
where $a\cdots c$, as elements of $\{\mathfrak G_\eta,\psi\}$, satisfy equations
with coefficients from $[\mathfrak G_\eta,\psi]$:
\[
\begin{gathered}
 a^\lambda+A_1a^{\lambda-1}+\cdots+A_\lambda=0,\\
 \vdots\\
 c^\mu+C_1c^{\mu-1}+\cdots+C_\mu=0.
\end{gathered}
\]
Let their roots be denoted by
$a,a'\cdots a^{(\lambda-1)}\cdots c,c'\cdots c^{(\mu-1)}$. If one now forms
the product over all combinations of the roots,
\[
 g(z)=f(z;a\cdots c)f(z;a'\cdots c')\cdots
      f(z;a^{(\lambda-1)}\cdots c^{(\mu-1)}),
\]
then $g(z)$ has unity as leading coefficient and elements of
$[\mathfrak G_\eta,\psi]$ as its remaining coefficients. Consequently $z$
belongs to $\{\mathfrak G_\eta,\psi\}$.\footnote{Thus, by defining
algebraic integrality by means of \emph{some} equation, the concept of
irreducibility with respect to a ground field becomes unnecessary for the proof
of I and V. Compare also § 2, where, only for the proof of IV, the auxiliary
proposition is needed which passes from the definition of algebraic integrality by
an arbitrary equation to the irreducible equation. Since this auxiliary proposition
no longer holds in general in the present case, IV has to be imposed on the
function $\psi$ as a special condition.}

Thus the domain $\{\mathfrak G_\eta,\psi\}$ has properties I and V; and, since
$\mathfrak G_\eta$ is a divisor of the domain, it also has II and III. Whether
IV is satisfied depends on the choice of $\psi$; for example, IV is not satisfied
for $\psi=\frac{1}{n\cdot\eta}$, since then $\frac1n$ belongs to the domain. It
is satisfied, however, by § 2 for $\psi=\frac1\eta$, or also for
$\psi=\frac{\eta}{n}$. For in the latter case, if a function from
$\{\mathfrak G_\eta,\psi\}$ represents an algebraic number, its value remains
unchanged for $\eta=0$; thereby it passes into a function from
$\mathfrak G_\eta$, and consequently into an algebraic integer.

Thus $\{\mathfrak G_\eta,\psi\}$ becomes a domain $\mathfrak H$ as soon as one
imposes on $\psi$ the condition that no algebraically fractional number enters the
domain through the algebraically integral adjunction.

Quite analogously one defines the rationally integral, respectively
algebraically integral, adjunction of an arbitrary system $\mathfrak S$ of
rational or algebraic functions of the $\eta$ to $\mathfrak G_\eta$. The
rationally integral adjunction leads to the integral domain
$[\mathfrak G_\eta,\mathfrak S]$, consisting of all integral polynomials in
finitely many elements from $\mathfrak G_\eta$ and $\mathfrak S$ at a time. The
algebraically integral adjunction gives the domain
$\{\mathfrak G_\eta,\mathfrak S\}$ consisting of all elements that depend
algebraically-integrally on $[\mathfrak G_\eta,\mathfrak S]$. Exactly as above,
one shows that the integral domain $\{\mathfrak G_\eta,\mathfrak S\}$ is
algebraically-integrally closed and also satisfies II and III. Thus, as above:

$\{\mathfrak G_\eta,\mathfrak S\}$ becomes a domain $\mathfrak H$ as soon as one
imposes on $\mathfrak S$ the condition that no algebraically fractional number
enters the domain through the algebraically integral adjunction --- that is, as
soon as $\mathfrak S$ becomes an admissible system.\footnote{More generally, one
shows analogously: if $\mathfrak T$ is any domain and $\{\mathfrak T\}$ denotes
the totality of elements depending algebraically-integrally on $\mathfrak T$,
then $\{\mathfrak T\}$ is algebraically-integrally closed.}

It is now important that the converse also holds, so that the domains
$\{\mathfrak G_\eta,\mathfrak S\}$ actually exhaust all domains $\mathfrak H$
from $\mathfrak M(H)$ as $\mathfrak S$ runs through all admissible systems.
Indeed, let $\mathfrak H$ be any domain from $\mathfrak M(H)$ and denote by
$\mathfrak L=\mathfrak H-\mathfrak G_\eta$ the system consisting of all functions
from $\mathfrak H$ not belonging to $\mathfrak G_\eta$. By V, $\mathfrak H$
contains the domain $\{\mathfrak G_\eta,\mathfrak L\}$ as a divisor; but
$\{\mathfrak G_\eta,\mathfrak L\}$ is also divisible by $\mathfrak H$, because
it contains $\mathfrak G_\eta$ and $\mathfrak L$, and since these have no common
element, it contains $\mathfrak H$ as well. Hence
$\mathfrak H=\{\mathfrak G_\eta,\mathfrak L\}$; and since IV is satisfied for
$\mathfrak H$, $\mathfrak L$ is of course an admissible system.\footnote{The
adjunction of a subsystem of $\mathfrak L$ can already suffice; for example, the
functional domain of § 2 arises by algebraically integral adjunction of all
integral rational functionals --- with the exception of the polynomials --- to
$\mathfrak G_\eta$.}

As $H$ runs through every possible algebraic basis, $\mathfrak M(H)$ runs, as
proved at the beginning, through the totality of all domains $\mathfrak H$.
Consequently the question of the most general domains $\mathfrak H$ from
algebraically-integral transcendental numbers is completely answered by the
results of this paragraph:
\begin{enumerate}
\item[a)] The intersection of all domains $\mathfrak H$ from $\mathfrak M(H)$ is
given by the Zermelo domain $\mathfrak G_\eta$, which itself is a domain
$\mathfrak H$; $\mathfrak M(H)$ is closed with respect to forming intersections.
\item[b)] Every domain $\mathfrak H$ from $\mathfrak M(H)$ arises by
algebraically integral adjunction of an admissible system $\mathfrak S$ to
$\mathfrak G_\eta$. As $H$ runs through every possible algebraic basis,
$\mathfrak M(H)$ runs through the totality of all domains $\mathfrak H$.\footnote{The
admissible systems $\mathfrak S$ can be found as follows. Put all algebraically
fractional functions of the $\eta$ under some well-ordering $\Omega$, and take
into $\mathfrak S$ all functions except those which, when algebraically-integrally
adjoined to $\mathfrak G_\eta$ and the previously admitted functions, generate an
algebraically fractional number. This procedure can be broken off at any arbitrary
place. As $\Omega$ runs through every well-ordering, and as the process is broken
off at every arbitrary place for each fixed $\Omega$, all admissible systems
$\mathfrak S$ are thereby exhausted.}
\end{enumerate}

\end{document}
